\documentclass[12pt]{article}

\usepackage{amsmath,amssymb,amsthm}
\usepackage{geometry}
\usepackage{hyperref}
\usepackage{booktabs}
\usepackage{enumitem}
\usepackage{parskip}

\geometry{letterpaper, margin=1in}

\newtheorem{theorem}{Theorem}[section]
\newtheorem{definition}[theorem]{Definition}
\newtheorem{proposition}[theorem]{Proposition}
\newtheorem{lemma}[theorem]{Lemma}
\newtheorem{conjecture}[theorem]{Conjecture}
\newtheorem{corollary}[theorem]{Corollary}
\newtheorem{remark}[theorem]{Remark}

% ── CPP macros ────────────────────────────────────────────────────────────────
\newcommand{\lP}{l_P}
\newcommand{\tP}{t_P}
\newcommand{\EP}{E_P}
\newcommand{\mP}{m_P}
\newcommand{\SSV}{\text{SSV}}
\newcommand{\PSR}{\text{PSR}}
\newcommand{\nuC}{\nu_C}
\newcommand{\rC}{r_C}
\newcommand{\rth}{r_{\rm th}}
\newcommand{\rin}{r_{\rm in}}
\newcommand{\rout}{r_{\rm out}}
\newcommand{\oin}{\omega_{\rm in}}
\newcommand{\oout}{\omega_{\rm out}}
\newcommand{\aBohr}{a_0}

\title{\textbf{Derivation of the Standing-Wave Sub-Harmonic Condition\\
for Captured Dipole Particle Orbits in Conscious Point Physics}\\[6pt]
{\large Spin Companion II ---
Companion Paper to ``Mechanistic Derivation of Relativistic Effects\\
via Space Stress Vector (SSV) in the Dipole Sea'' (Version~16)}}

\author{Thomas Lee Abshier, ND\\
Hyperphysics Institute\\
\texttt{https://hyperphysics.com}\\
\texttt{drthomas007@protonmail.com}}

\date{16 March 2026 --- Version~2}

\begin{document}
\maketitle

\noindent\textbf{Keywords:} Conscious Point Physics, emergent spin,
standing wave, sub-harmonic, ZBW cloud, orbital geometry, fine structure
constant, Bohr radius, Compton wavelength.

%======================================================================
\section*{Plain Language Summary}
%======================================================================

The Spin~I companion paper derived the electron spin quantum number
$\hbar/2$ from the orbital angular momentum of a captured Dipole
Particle, given one geometric condition: the outer orbital radius is
exactly twice the inner orbital radius ($\rout = 2\rin$).  That
condition was physically motivated but not derived.  This paper
derives it.

The Zitterbewegung (ZBW) cloud surrounding every massive particle is
a standing wave oscillating between a fixed inner wall and a free outer
boundary.  Standing waves in such a cavity have a ladder of allowed
vibration patterns: the first pattern has no internal structure, the
second has one internal peak and one internal zero, the third has more.
A Dipole Particle captured by this cloud must sit at an internal peak
(where the field is maximum) for its positive CP and at an internal
zero (where the net force is zero) for its negative CP.  The second
vibration pattern is the first one that provides both.  In that pattern,
the internal peak sits at exactly one-third of the cavity size, and the
internal zero sits at exactly two-thirds.  The ratio is exactly 2:1.

This is the geometric origin of $\rout = 2\rin$.  Combined with the
Spin~I result, it completes the first-principles derivation of spin
$1/2$ from the CPP lattice geometry.

The absolute size of the orbit is not set by the ZBW cloud size
but by the Coulomb force balance --- this is the Bohr radius divided
by a geometric factor, as derived in Spin~I.  The connection between
the ZBW cloud scale and the Bohr scale involves the fine structure
constant $\alpha \approx 1/137$ in a specific way that is identified
here and deferred for further investigation.

%======================================================================
\begin{abstract}
The Spin~I companion established that the electron spin quantum
number $L = \hbar/2$ follows from the captured Dipole Particle orbital
geometry given the condition $\rout = 2\rin$.  We derive this condition
from the standing-wave mode structure of the ZBW polarization cloud.
The ZBW cloud constitutes a radial standing-wave resonator with a node
at $r = 0$ (CP Exclusion Rule) and an antinode at the thermal boundary
$\rth = \hbar/(2m_e c)$.  The allowed modes are:
\begin{equation}
  \psi_n(r) \;=\; \sin(k_n r),
  \qquad
  k_n \;=\; \frac{(2n-1)\pi}{2\rth},
  \quad n = 1,\,2,\,3,\ldots
  \label{eq:modes_abstract}
\end{equation}
The fundamental mode ($n=1$) has no interior structure.  The first
sub-harmonic ($n=2$) has an interior antinode at $\rth/3$ and an
interior node at $2\rth/3$.  The captured DP anchors its inner $+$CP
at the interior antinode and its outer $-$CP at the interior node,
giving:
\begin{equation}
  \frac{\rout}{\rin} \;=\; \frac{2\rth/3}{\rth/3} \;=\; 2.
  \label{eq:ratio_abstract}
\end{equation}
This ratio is independent of all physical constants and depends only
on the integer arithmetic of the trigonometric zeros of the $n=2$
mode.
This is exact and independent of all physical constants.  Mode 2 is
the minimum-energy mode that simultaneously provides an interior
antinode and an interior node.  The orbital scale differs from the ZBW
cloud scale by the factor $6/[\alpha\cdot 4(1+\sqrt{2})^2]$ where
$\alpha$ is the fine structure constant, connecting the Compton scale
to the Bohr scale.  The derivation of this scale connection from the
600-cell Voronoi eigenvalue spectrum is deferred to Spin~III.
\end{abstract}

%======================================================================
\section{Introduction}
\label{sec:intro}
%======================================================================

The Spin~I companion~\cite{sI} derived the electron spin quantum number
from first principles in CPP.  The central result was
Theorem~5.1 of that paper:
\begin{equation}
  L \;=\; m_e\,\oout\,\rin^2\,(2\sqrt{2}+4) \;=\; \frac{\hbar}{2}
  \qquad\text{when}\qquad
  \rin \;=\; \frac{\aBohr}{4(1+\sqrt{2})^2},
  \label{eq:spinI_result}
\end{equation}
given the geometric condition $\rout = 2\rin$.  That condition was
stated as Open Problem~1 of Spin~I: it is physically motivated by
the ZBW standing-wave structure, but the derivation from the 600-cell
geometry was deferred.

The present paper closes this gap in the sub-harmonic regime.  We
show that the condition $\rout = 2\rin$ follows rigorously from the
mode structure of the ZBW radial standing wave, with no additional
assumptions.  The argument uses only:
\begin{enumerate}[noitemsep]
  \item The ZBW cloud boundary conditions from the ZBW Mass
    companion~\cite{c4}.
  \item The standard mathematics of standing waves in a
    node-antinode resonator.
  \item The physical requirement that the captured DP anchors
    its inner $+$CP at a field maximum and its outer $-$CP at
    a field zero.
\end{enumerate}

The connection between the ZBW cloud scale ($\rth$) and the Bohr
orbital scale ($\aBohr$) is identified as a separate problem involving
the fine structure constant $\alpha$, and is deferred to Spin~III.

%======================================================================
\section{The ZBW Radial Standing Wave}
\label{sec:ZBW_wave}
%======================================================================

\subsection{Boundary conditions from the ZBW Mass companion}

The ZBW Mass companion~\cite{c4} established that every massive
particle is surrounded by a compressive eDP polarization cloud.
The cloud is a radial standing wave with two boundary conditions:

\begin{definition}[ZBW resonator boundary conditions]
\label{def:BC}
\quad
\begin{enumerate}[noitemsep]
  \item \emph{Inner boundary (node):} $\psi(0) = 0$.  The CP
    Exclusion Rule~\cite{c1} prevents any CP displacement at
    $r = 0$; this is a displacement node.
  \item \emph{Outer boundary (antinode):} $\psi(\rth) = $ maximum.
    The thermal boundary $\rth = \hbar/(2m_e c) = \rC/2$ is a
    free boundary where the restoring force vanishes; this is a
    displacement antinode.
\end{enumerate}
\end{definition}

These are the boundary conditions of an open-closed resonator
(analogous to an organ pipe closed at one end and open at the other,
or a string fixed at one end and free at the other).

\subsection{Allowed mode spectrum}

\begin{proposition}[ZBW mode spectrum]
\label{prop:modes}
The allowed radial modes satisfying Definition~\ref{def:BC} are:
\begin{equation}
  \psi_n(r) \;=\; A\sin(k_n r),
  \qquad
  k_n \;=\; \frac{(2n-1)\pi}{2\rth},
  \quad n = 1,\,2,\,3,\ldots
  \label{eq:modes}
\end{equation}
with wavelengths $\lambda_n = 4\rth/(2n-1)$.
\end{proposition}

\begin{proof}
The condition $\psi(0) = 0$ requires $\psi \propto \sin(kr)$.
The condition $\psi'(\rth) = 0$ (maximum at outer boundary) requires
$k\rth = \pi/2 + m\pi$ for non-negative integer $m$.  Setting
$m = n-1$ gives $k_n = (2n-1)\pi/(2\rth)$.
\end{proof}

The mode energies scale as $k_n^2 \propto (2n-1)^2$:
\begin{center}
\renewcommand{\arraystretch}{1.3}
\begin{tabular}{@{}ccccc@{}}
\toprule
Mode $n$ & $k_n$ (units of $\pi/2\rth$) & $\lambda_n$ &
Interior nodes & Interior antinodes \\
\midrule
1 (fundamental)   & 1 & $4\rth$         & none            & none \\
2 (1st sub-harmonic) & 3 & $\tfrac{4}{3}\rth$ & $\tfrac{2}{3}\rth$ & $\tfrac{1}{3}\rth$ \\
3 (2nd sub-harmonic) & 5 & $\tfrac{4}{5}\rth$ & $\tfrac{2}{5}\rth$, $\tfrac{4}{5}\rth$ & $\tfrac{3}{5}\rth$ \\
\bottomrule
\end{tabular}
\end{center}

The fundamental mode has no interior structure --- it reaches its
maximum only at the outer boundary $\rth$.  Only modes $n \geq 2$
have both interior nodes and interior antinodes.

%======================================================================
\section{The Captured DP Anchoring Conditions}
\label{sec:anchoring}
%======================================================================

The captured DP is a bound pair of CPs (inner $+$eCP at $\rin$, outer
$-$eCP at $\rout$) in stable circular orbit about the central $-$eCP.
Stable anchoring requires:

\begin{definition}[DP anchoring conditions]
\label{def:anchor}
\quad
\begin{enumerate}[noitemsep]
  \item \emph{Inner $+$CP at field maximum (antinode):} The inner
    $+$CP is attracted to the central $-$eCP.  Stable anchoring
    requires it to sit at a displacement maximum of the ZBW field
    --- an interior antinode --- where the restoring force is zero
    and the field-mediated Coulomb potential is at its local
    extremum.
  \item \emph{Outer $-$CP at field zero (node):} The outer $-$eCP
    is repelled from the central $-$eCP.  Stable anchoring requires
    it to sit at a displacement node where the net SSV force
    averages to zero over a ZBW cycle, preventing cumulative drift.
\end{enumerate}
\end{definition}

These conditions identify the relevant sub-harmonic mode:

\begin{lemma}[Mode selection]
\label{lem:mode_select}
The minimum-energy mode satisfying both anchoring conditions of
Definition~\ref{def:anchor} is $n = 2$.
\end{lemma}

\begin{proof}
Mode $n=1$ has no interior node or antinode --- it cannot satisfy
either anchoring condition.

Mode $n=2$ has exactly one interior antinode (at $\rth/3$) and
one interior node (at $2\rth/3$), satisfying both conditions
simultaneously.  Its energy $\propto k_2^2 = 9(\pi/2\rth)^2$.

Mode $n=3$ and higher also satisfy both conditions but at higher
energy $\propto (2n-1)^2 > 9$.  By minimum energy, the system
selects $n=2$.
\end{proof}

%======================================================================
\section{Derivation of $\rout = 2\rin$}
\label{sec:derivation}
%======================================================================

\begin{theorem}[$\rout = 2\rin$ from ZBW sub-harmonic]
\label{thm:ratio}
The DP anchoring conditions (Definition~\ref{def:anchor}) applied to
the minimum-energy eligible mode (Lemma~\ref{lem:mode_select}) give
exactly:
\begin{equation}
  \rout \;=\; 2\rin.
  \label{eq:ratio}
\end{equation}
\end{theorem}

\begin{proof}
For Mode $n=2$, the wave function is:
\begin{equation}
  \psi_2(r) \;=\; A\sin\!\left(\frac{3\pi r}{2\rth}\right).
  \label{eq:psi2}
\end{equation}
The interior antinode satisfies $\psi_2'(r) = 0$:
\begin{equation}
  \frac{3\pi}{2\rth}\cos\!\left(\frac{3\pi r}{2\rth}\right) = 0
  \;\Rightarrow\; \frac{3\pi r}{2\rth} = \frac{\pi}{2}
  \;\Rightarrow\; r = \frac{\rth}{3}.
  \label{eq:antinode_pos}
\end{equation}
The interior node satisfies $\psi_2(r) = 0$:
\begin{equation}
  \sin\!\left(\frac{3\pi r}{2\rth}\right) = 0
  \;\Rightarrow\; \frac{3\pi r}{2\rth} = \pi
  \;\Rightarrow\; r = \frac{2\rth}{3}.
  \label{eq:node_pos}
\end{equation}
By Definition~\ref{def:anchor}, the inner $+$CP anchors at the
antinode: $\rin = \rth/3$.  The outer $-$CP anchors at the node:
$\rout = 2\rth/3$.  Therefore:
\begin{equation}
  \frac{\rout}{\rin} \;=\; \frac{2\rth/3}{\rth/3} \;=\; 2.
  \label{eq:ratio_proof}
\end{equation}
\end{proof}

\begin{remark}[Independence from physical constants]
The ratio $\rout/\rin = 2$ depends only on the positions of the
interior node and antinode in Mode~2, which are fixed by the
integer arithmetic of the trigonometric zeros.  It is independent
of $\rth$, $m_e$, $e$, $\hbar$, $k_e$, and all other physical
constants.  This is the geometric origin of the $2\sqrt{2}$ angular
frequency ratio in Spin~I.
\end{remark}

\begin{remark}[Connection to the $2\sqrt{2}$ ratio]
The $2\sqrt{2}$ angular-frequency ratio of Spin~I follows immediately
from $\rout = 2\rin$ via the $1/r^2$ Coulomb force law
(Spin~I, Proposition~3.1):
\begin{equation}
  \frac{\oin}{\oout} \;=\; \left(\frac{\rout}{\rin}\right)^{3/2}
  \;=\; 2^{3/2} \;=\; 2\sqrt{2}.
\end{equation}
Theorem~\ref{thm:ratio} of this paper is therefore the geometric
foundation for the entire Spin~I derivation.
\end{remark}

%======================================================================
\section{Numerical Verification}
\label{sec:numerical}
%======================================================================

The mode-2 wave function is verified at all critical points:

\begin{center}
\renewcommand{\arraystretch}{1.4}
\begin{tabular}{@{}lccl@{}}
\toprule
Point & $r$ & $\psi_2(r)$ & Identification \\
\midrule
Inner boundary  & $0$           & $0$   & Node (CP Exclusion) \\
Inner $+$CP     & $\rth/3$      & $+1$  & Interior antinode ($\rin$) \\
Outer $-$CP     & $2\rth/3$     & $0$   & Interior node ($\rout$) \\
Outer boundary  & $\rth$        & $-1$  & Boundary antinode \\
\midrule
Ratio $\rout/\rin$ & --- & --- & $2.000\,000$ (exact) \\
\bottomrule
\end{tabular}
\end{center}

Numerical values:
\begin{align}
  \rth   &\;=\; \frac{\hbar}{2m_e c}
  \;=\; 1.9308\times 10^{-13}\ \text{m}, \notag\\
  \rin^{\rm (ZBW)} &\;=\; \frac{\rth}{3}
  \;=\; 6.436\times 10^{-14}\ \text{m}, \notag\\
  \rout^{\rm (ZBW)} &\;=\; \frac{2\rth}{3}
  \;=\; 1.287\times 10^{-13}\ \text{m}.
  \label{eq:numerical}
\end{align}

%======================================================================
\section{The Scale Connection: ZBW Cloud vs.\ Orbital Scale}
\label{sec:scale}
%======================================================================

The ZBW sub-harmonic fixes the \emph{ratio} $\rout = 2\rin$ but not
the absolute orbital scale.  The orbital scale is determined by the
Coulomb force balance and spin quantization condition of Spin~I:
\begin{equation}
  \rin^{\rm (orbit)} \;=\; \frac{\aBohr}{4(1+\sqrt{2})^2}
  \;=\; 2.270\times 10^{-12}\ \text{m}.
  \label{eq:orbital_scale}
\end{equation}
The ratio of orbital scale to ZBW cloud scale is:
\begin{equation}
  \frac{\rin^{\rm (orbit)}}{\rin^{\rm (ZBW)}}
  \;=\; \frac{\aBohr/[4(1+\sqrt{2})^2]}{\rth/3}
  \;=\; \frac{3\aBohr}{4(1+\sqrt{2})^2\,\rth}
  \;=\; \frac{6}{\alpha\cdot 4(1+\sqrt{2})^2}
  \;\approx\; 35.27,
  \label{eq:scale_ratio}
\end{equation}
where we used $\aBohr = \rth\cdot 2/\alpha$ (since $\aBohr = \rC/\alpha$
and $\rth = \rC/2$) and $\alpha \approx 1/137.036$ is the fine
structure constant.

\begin{remark}[Physical interpretation of the scale ratio]
The factor $\alpha \approx 1/137$ that enters Eq.~\eqref{eq:scale_ratio}
is the ratio of the electromagnetic interaction energy to the rest-mass
energy at the Compton scale.  The Bohr radius exceeds the Compton
radius by exactly $1/\alpha$, because the Coulomb binding energy at
the Bohr radius equals the electron rest mass times $\alpha^2$.

The scale ratio~\eqref{eq:scale_ratio} therefore means: the orbital
DP sits at a Coulomb-bound scale ($\sim\aBohr$) that is parametrically
larger than the ZBW cloud scale ($\sim\rth$) by the standard
Compton-to-Bohr ratio, divided by the geometric factor $4(1+\sqrt{2})^2$.
The sub-harmonic fixes the ratio of the two orbital radii within the
orbit; the Coulomb interaction determines where within the potential
well the orbit sits.
\end{remark}

The two derivations are therefore complementary and non-redundant:
\begin{itemize}[noitemsep]
  \item \textbf{Spin~II (this paper):} ZBW sub-harmonic $\to$
    $\rout/\rin = 2$ (ratio only).
  \item \textbf{Spin~I:} Coulomb force balance + $L = \hbar/2$
    $\to$ $\rin = \aBohr/[4(1+\sqrt{2})^2]$ (scale only).
  \item \textbf{Together:} orbital geometry is fully determined.
\end{itemize}

%======================================================================
\section{Consistency with Spin~I}
\label{sec:consistency_spinI}
%======================================================================

Theorem~\ref{thm:ratio} of this paper provides the geometric
foundation for Proposition~3.1 and Theorem~5.1 of Spin~I:

\begin{enumerate}[noitemsep]
  \item $\rout = 2\rin$ (Theorem~\ref{thm:ratio} here) is the input
    to the Spin~I force balance.
  \item Given this ratio, the Spin~I force balance gives
    $\oin/\oout = 2\sqrt{2}$ exactly (Spin~I, Proposition~3.1).
  \item Given this ratio and the Coulomb scale, the Spin~I angular
    momentum calculation gives $L = \hbar/2$ exactly (Spin~I,
    Theorem~5.1).
\end{enumerate}

The two papers together constitute a complete derivation of the
electron spin quantum number from CPP geometry and the Coulomb
interaction, with no free parameters.

%======================================================================
\section{Consistency with Previous Companions}
\label{sec:consistency}
%======================================================================

\begin{itemize}[noitemsep]
  \item The \emph{ZBW standing wave} and thermal boundary
    $\rth = \rC/2$ are from the ZBW Mass companion~\cite{c4},
    Proposition~4.1.
  \item The \emph{CP Exclusion Rule} providing the inner node at
    $r = 0$ is from the Absolute Moment companion~\cite{c1}, \S2.
  \item The \emph{node-antinode resonator} is the continuum limit
    of the 600-cell lattice standing wave; the exact lattice
    version uses the Voronoi eigenvalue spectrum (deferred).
  \item The \emph{fine structure constant} $\alpha$ connecting
    the two scales enters through $\aBohr = \rC/\alpha$, which
    follows from the Stiffness~$C$ companion~\cite{c2}, \S4.3.
  \item No element contradicts any previous companion.
\end{itemize}

%======================================================================
\section{Open Problems}
\label{sec:open}
%======================================================================

\textbf{(1) Spin~III --- Full 600-cell Voronoi eigenvalue spectrum.}
The node-antinode boundary conditions and mode spectrum used here are
the continuum (spherical) limit.  The exact CPP statement is that the
ZBW standing wave lives on the 600-cell lattice, whose Voronoi cells
are 24-cells.  The eigenvalue spectrum of the Laplacian on the 24-cell
Voronoi decomposition will differ from the spherical limit by terms of
order $(l_P/\rth)^2 \sim 10^{-44}$ --- negligible at any measurable
scale, but required for a fully first-principles derivation.  This
is the subject of the next companion: Spin~III.

\textbf{(2) Mode selection from first principles.}
Lemma~\ref{lem:mode_select} identifies Mode~2 as the
minimum-energy mode satisfying the anchoring conditions.  A more
rigorous argument would derive the selection from the energy landscape
of the full DP--ZBW interaction, including the DP binding energy and
the thermal disruption term.

\textbf{(3) The scale ratio and $\alpha$.}
The factor $6/[\alpha\cdot 4(1+\sqrt{2})^2]$ in
Eq.~\eqref{eq:scale_ratio} connects the ZBW cloud scale to the Bohr
orbital scale through the fine structure constant.  Whether this
ratio has a deeper geometric meaning --- perhaps from the 600-cell
coordination number or the $H_4$ Coxeter symmetry --- is an open
question.

\textbf{(4) Higher modes and excited states.}
Modes $n = 3, 4, \ldots$ give higher sub-harmonics with multiple
interior nodes and antinodes.  These correspond to excited orbital
configurations of the captured DP, with angular momenta $L > \hbar/2$.
Whether these map to the excited lepton states (muon, tau) or to
excited orbital shells is under investigation.

%======================================================================
\section{Conclusion}
\label{sec:conclusion}
%======================================================================

\begin{enumerate}
  \item \textbf{ZBW mode spectrum} (rigorous,
    Proposition~\ref{prop:modes}): the ZBW radial standing wave
    has a node-antinode resonator spectrum
    $k_n = (2n-1)\pi/(2\rth)$, with Mode~1 having no interior
    structure and Mode~2 being the first to have both an interior
    antinode and an interior node.

  \item \textbf{Mode selection} (Lemma~\ref{lem:mode_select}):
    Mode~2 is the minimum-energy mode simultaneously satisfying
    both DP anchoring conditions.

  \item \textbf{$\rout = 2\rin$} (rigorous, Theorem~\ref{thm:ratio}):
    the interior antinode of Mode~2 is at $\rth/3$ and the interior
    node is at $2\rth/3$.  Their ratio is exactly 2, independent of
    all physical constants.

  \item \textbf{Scale connection} (identified, deferred): the orbital
    scale ($\aBohr/[4(1+\sqrt{2})^2]$) exceeds the ZBW cloud scale
    ($\rth/3$) by the factor $6/[\alpha\cdot 4(1+\sqrt{2})^2] \approx
    35.27$, connecting the Compton scale to the Bohr scale through
    $\alpha$.

  \item \textbf{Completion of spin derivation}: Theorem~\ref{thm:ratio}
    of this paper combined with Theorem~5.1 of Spin~I constitutes a
    complete first-principles derivation of the electron spin quantum
    number $\hbar/2$ from the CPP geometry, the $1/r^2$ Coulomb law,
    and the ZBW standing-wave structure.  The one remaining step
    (600-cell Voronoi eigenvalue proof) is deferred to Spin~III.
\end{enumerate}

%======================================================================
\section*{Acknowledgments}
%======================================================================

The identification of the ZBW sub-harmonic as the geometric origin of
$\rout = 2\rin$ was developed in collaboration with Grok (xAI) and
Claude (Anthropic).  The numerical verification and mode analysis were
carried out by Claude.  The author thanks both for contributions to the
CPP spin series.

%======================================================================
\begin{thebibliography}{99}

\bibitem{sI}
T.~L. Abshier,
``Emergent Spin-$\tfrac{1}{2}$ from Captured Dipole Particle Orbital
Geometry in Conscious Point Physics,''
Version~2, 2026 (Spin companion~I).

\bibitem{v16}
T.~L. Abshier,
``Mechanistic Derivation of Relativistic Effects via Space Stress
Vector (SSV) in the Dipole Sea,''
Version~16, 2026 (main paper).

\bibitem{c1}
T.~L. Abshier,
``The Absolute Moment Postulate,''
Version~3, 2026 (companion~1).

\bibitem{c2}
T.~L. Abshier,
``Microscopic Origin of the Dipole Sea Stiffness~$C$,''
Version~2, 2026 (companion~2).

\bibitem{c3}
T.~L. Abshier,
``Quantum Probability and the Classical Transition in CPP,''
Version~2, 2026 (companion~3).

\bibitem{c4}
T.~L. Abshier,
``Inertial Mass from Zitterbewegung,''
Version~2, 2026 (companion~4).

\bibitem{c5}
T.~L. Abshier,
``Newtonian Gravity from SSV Shell Broadcast,''
Version~2, 2026 (companion~5).

\bibitem{c6}
T.~L. Abshier,
``Dipole Chain Patterns as the Substrate of Mass and
Electromagnetic Radiation,''
Version~2, 2026 (companion~6).

\bibitem{c7}
T.~L. Abshier,
``Weak-Field General Relativity from the Lattice State Packet,''
Version~3, 2026 (companion~7).

\bibitem{dirac1928}
P.~A.~M. Dirac,
``The Quantum Theory of the Electron,''
\textit{Proc.\ R.\ Soc.\ London A}, \textbf{117}, 610--624, 1928.

\bibitem{codata}
NIST, CODATA 2018 Recommended Values of the Fundamental Physical
Constants,
\url{https://physics.nist.gov/cuu/Constants/}, 2018.

\end{thebibliography}

\end{document}
